\begin{theorem}[Proved: protected first-stage cycles and a split-monic quotient]\label{C:protected-cycles}
In any repaired AKR Construction C.1, let $P_+$ and $P_-$ be the fresh
first-stage quasi-filter targets for $(v_1,v_2)$ and $(v_2,v_1)$.
Write $s_+,s_-$ for their respective numbers of quasi-filter maps.
In the binary graph $G_{P_\pm}$, let $S_\pm$ be the $s_\pm$ vertices
of its unique circular component, and choose its integral cycle
$\kappa_\pm$ with augmentation one. For $X\in K^*$ define
\[
 F(X)=\mathbb Z[K^*(P_+,X)]\oplus\mathbb Z[K^*(P_-,X)],
 \qquad
 \iota_X(b)=b\kappa_\pm\in K(X)
 \quad(b:P_\pm\to X).
\]
Then $\iota:F\to K$ is a natural monomorphism. Its quotient $M=K/F$
is pointwise a finitely generated free abelian group, and every arrow
of $K^*$ induces a split injection on $M$ with free abelian cokernel.
The splitting is an existence statement for each individual arrow;
no compatible family of splittings is asserted.

More precisely, if $X\in K_m$, $m\geq1$, $g:Y\to X$ is any arrow,
and $h_i:P_{\epsilon_i}\to X$ is a finite family of distinct typed
arrows, then:
\begin{enumerate}
\item $|h_i(S_{\epsilon_i})\cap g(Y)|\geq3$ implies that $h_i$
factors through $g$;
\item if every displayed intersection has at most two points, then
\[
 \left|\bigcup_i h_i(S_{\epsilon_i})\setminus g(Y)\right|
 \geq\frac{a_m}{a_1}\sum_i(s_{\epsilon_i}-2).
\]
\end{enumerate}
These assertions hold for mixed families with either source $P_+$ or
$P_-$.
\end{theorem}
\begin{proof}
We first verify the marked-set versions of the source induction.
At stage one, a map from $P_+$ or $P_-$ to an old object is its identity,
if the target is that same object, and does not otherwise exist.
Every proper incoming map to $P_\pm$ has source $V$ or $W$. Its image
meets the marked circle in at most two points: a $W$-map supplies one
edge of that circle and its third point lies outside it. Therefore the
first assertion holds at stage one, and the second has bound
$s_\pm-2$ for the only possible nonempty family. There are no arrows
between $P_+$ and $P_-$, and the only endomorphisms of either are
identities; this remains so because all stages are full extensions.

We give the quasi-filter step explicitly. Write the old pair as
$\phi_1,\phi_2:R\rightrightarrows T$ and its new maps as
$\alpha_j:T\to Z$, with $j\in\mathbb Z/s$. An identity comparison
map $g=1_Z$ makes the first assertion trivial and the second vacuous.
Otherwise write $g=\alpha_0g^*$ after a cyclic renumbering. Every marked
arrow into $Z$ factors as $h=\alpha_jh^*$.

Suppose first that $h(S_\epsilon)$ meets $g(Y)$ in at least three
points. If $1\leq j\leq s-2n$, (10) gives
$\alpha_j(T)\cap\alpha_0(T)\subseteq
\alpha_j(T)\cap\alpha_{j-1}(T)$. If $s-2n<j<s$, move instead toward
$j+1$; its forward distance to $0$ is less than $2n\leq s-2n$, so the
other containment in (10) applies. In either case choose the adjacent
identity
\[
 \alpha_j\phi_p=\alpha_{j'}\phi_q,
 \qquad j'=j-1\ \text{or}\ j+1,
 \quad p,q\in\{1,2\}.
\]
By (11) the intersection of the two copies is exactly
$\alpha_j\phi_p(R)$. Injectivity of $\alpha_j$ therefore gives
$|h^*(S_\epsilon)\cap\phi_p(R)|\geq3$. The old first assertion
supplies $l:P_\epsilon\to R$ with $h^*=\phi_pl$. Hence
$h=\alpha_{j'}\phi_ql$: the whole typed arrow, not merely its marked
set, has moved to the adjacent copy. Its image and its marked image
are unchanged. If we started with $j\leq s-2n$, repeat downward to
$0$; otherwise repeat upward through $s-1$ to $0$. Each relevant
distance decreases, so (10) remains applicable. Once
$h=\alpha_0h_0^*$, injectivity of $\alpha_0$ and the old first
assertion applied to $h_0^*,g^*$ show that $h$ factors through $g$.
This proves the new first assertion without any use of the number of
points of $P_\epsilon$.

For the second assertion let the family be $h_1,\ldots,h_r$, give
member $i$ the positive weight $w_i=s_{\epsilon_i}-2$, and put
$\beta_m=a_m/a_1$. Define disjoint sets
\[
 H_j=\{i:h_i\text{ factors through }\alpha_j\}
          \setminus\bigcup_{t<j}H_t,
 \qquad 0\leq j<s.
\]
They partition the family. Partition the indices $0,\ldots,s-1$ into
$s/(3n)$ consecutive blocks of length $3n$. Choose a block
$B=\{j_0,\ldots,j_0+3n-1\}$ minimizing
$\sum_{j\in B}\sum_{i\in H_j}w_i$, and discard its assigned family
members. For the retained index set
$L=\{1,\ldots,r\}\setminus\bigcup_{j\in B}H_j$ we have
\[
 \sum_{i\notin L}w_i\leq\frac{3n}{s}\sum_{i=1}^r w_i,
 \qquad
 \sum_{i\in L}w_i\geq\frac{a_{m+1}}{a_m}\sum_{i=1}^r w_i.
\]
Now order the remaining copy indices by
\[
 0,1,\ldots,j_0-1,\quad s-1,s-2,\ldots,j_0+3n,
\]
omitting the first arm if $j_0=0$. Assign each retained member to
the first available copy in this order. Denote by $L_j$ the family
assigned to $j$, write $h_i=\alpha_jh_i^*$ for $i\in L_j$, and put
$D_j=\bigcup_{i\in L_j}h_i(S_{\epsilon_i})$.

If the first index is $0$, the old second assertion, applied to
$g^*$ and the $h_i^*$ assigned there, yields
\[
 |D_0\setminus g(Y)|\geq\beta_m\sum_{i\in L_0}w_i.
\]
For every later index $j$ in the first arm, all intersections with
earlier copies and with $g(Y)$ lie, by (10), in
$\alpha_j(T)\cap\alpha_{j-1}(T)$. Indeed all backward distances
are at most $j_0-1<s-2n$. For an index $j$ in the second arm, every
earlier index, and also $0$, is reached in the forward cyclic
direction in at most $s-3n$ steps, so those intersections lie in
$\alpha_j(T)\cap\alpha_{j+1}(T)$. This includes the first index
$s-1$ when $j_0=0$.

Let $j'$ be this adjacent comparison index. No member of $L_j$
factors through $\alpha_{j'}$. Usually $j'$ is earlier in the sweep,
so this follows from the assignment rule. In the only exceptional
case $j=s-1,j_0=0$, a factorization through $\alpha_0$ would put that
member in $H_0$, hence in the discarded block, again impossible.
Choose $\alpha_j\phi_p=\alpha_{j'}\phi_q$ as above. If
$h_i^*(S_{\epsilon_i})$ met $\phi_p(R)$ in three points, the old
first assertion would factor $h_i^*$ through $\phi_p$ and hence
$h_i$ through $\alpha_{j'}$, a contradiction. Therefore the old
second assertion applies to the family $h_i^*$ and comparison arrow
$\phi_p$, and (11) gives
\[
 \begin{split}
 \left|D_j\setminus\left(g(Y)\cup
               \bigcup_{t\text{ earlier}}D_t\right)\right|
 &\geq\left|D_j\setminus\alpha_j\phi_p(R)\right|\\
 &=\left|\bigcup_{i\in L_j}h_i^*(S_{\epsilon_i})
                         \setminus\phi_p(R)\right|\\
 &\geq\beta_m\sum_{i\in L_j}w_i.
 \end{split}
\]
Summing these disjoint increments over the sweep and using the
retained-weight estimate proves
\[
 \left|\bigcup_i h_i(S_{\epsilon_i})\setminus g(Y)\right|
 \geq\beta_m\sum_{i\in L}w_i
 \geq\beta_{m+1}\sum_{i=1}^r w_i.
\]
Each factorization preserves the source and its weight; injectivity
of a copy map preserves distinctness of the typed arrows. Thus the
argument is valid for mixed protected source types. It uses only
(10), (11), monicity, the old threshold-three assertion, and the old
weighted bound. The original special cardinality $|W|=3$ is not
used in this inductive step.

At a disjoint-sum target, the two summands have disjoint point sets.
If the comparison arrow is the identity of the new target, the
factorization assertion is immediate and the nonfactorization bound
is vacuous. Otherwise every nonempty marked overlap places the marked
arrow and the comparison arrow in the same summand; the old assertion
then gives the required factorization.
Apply the old bound in the comparison summand. In the other summand
one may use the following unconditional bound. Choose one marked
arrow $h_r$ as comparison. Every other member has at most two marked
points in its full image, since otherwise it would factor through
$h_r$. Such a factorization between the two fixed source objects is
impossible unless both sources and arrows coincide. Applying the old
second assertion to the other members and then adding all
$s_{\epsilon_r}$ points of $h_r(S_{\epsilon_r})$ gives at least
$(a_m/a_1)\sum_iw_i$, because
$(a_m/a_1)(s_{\epsilon_r}-2)\leq s_{\epsilon_r}$. This proves both
marked assertions for all stages.

Every admissible first-stage length satisfies $s_\pm\geq45$.
Indeed its active path has two vertices, divisibility
$3n\mid s=4k+1$ forces $n$ odd, hence $n\geq3$, and
$s/(3n)\geq(1-a_1/2)^{-1}>4$ is an integer. Also
$a_m/a_1>3/4$. Consequently
\[
 \frac{a_m}{a_1}(s_\pm-2)>\frac{s_\pm}{2}.
\]
The unconditional bound from the preceding paragraph shows that every
nonempty finite family of marked circles has a vertex belonging to
exactly one of them: otherwise the union would have at most half the
total number of marked incidences. At such a vertex, an incident edge
of that circle belongs to no other circle in the family. Its
coefficient detects the coefficient of that circle in a proposed
linear relation. This proves injectivity of $\iota_X$, over
$\mathbb Z$ and indeed with arbitrary abelian-group coefficients.
Naturality follows from composition of arrows.

For the quotient assertion, attach a two-cell to $G_X$ along each
marked embedded circle $b\kappa_\pm$ and denote the resulting finite
CW complex by $C_X$. The preceding private-vertex argument provides
a boundary edge of one two-cell that belongs to no other two-cell.
Collapsing that cell across that edge and repeating removes all
two-cells. Thus $C_X$ is homotopy equivalent to a graph, and
\[
 H_1(C_X;\mathbb Z)=K(X)/F(X)=M(X)
\]
is free abelian.

An arrow $g:Y\to X$ embeds $C_Y$ cellularly into $C_X$: it is
injective on vertices, edges and typed protected cells. The cells not
in its image correspond exactly to protected arrows that do not
factor through $g$. For any nonempty family of those cells the marked
assertions give
\[
 \left|\bigcup_i h_i(S_{\epsilon_i})\setminus g(Y)\right|
 >\frac12\sum_i s_{\epsilon_i}.
\]
Hence there is a private marked vertex outside $g(Y)$. A boundary
edge incident to it is free among the new two-cells and is not an
edge of $g(C_Y)$. We can therefore perform all collapses relative
to $g(C_Y)$. The resulting relative CW complex has only additional
zero- and one-cells. Adding such cells injects each old first homology
group and has free abelian cokernel, as follows directly by extending
a maximal forest. Thus $M(g)$ is injective with free cokernel and,
over $\mathbb Z$, is split injective.
\end{proof}

\paragraph{The remaining colimit obstruction after this theorem.}
The theorem is pointwise and relative to individual arrows. It does
not assert injectivity after taking a colimit. The long exact sequence
of $0\to F\to K\to M\to0$ contains
\[
 T_1(M)\longrightarrow\mathbb Z^2
 \longrightarrow\operatorname{colim}_{K^*}K.
\]
The two standard basis vectors are the two protected cycles. Their
augmentation difference $(1,-1)$ is the first-stage two-dimensional
class. The theorem has not shown that this difference lies outside
the image of $T_1(M)$.

